% Moins taxer les ventes, mieux taxer le sol (rapport Sorbon, partie III, mesure 4)
% Compilation depuis fiscalite/foncier/ : pdflatex note-foncier.tex (deux passes)
% Chiffres : python3 simulation.py -> sorties/simulation.txt ; état des lieux : notes/inventaire.md ;
% littérature : notes/droits_mutation.md, notes/taxe_fonciere_split_rate.md, ../lvt/notes/
% Figures nues : FIG_NU=1 python3 foncier/figures_foncier.py (depuis fiscalite/)
\documentclass[11pt,a4paper]{article}
% Préambule commun au brief et à la revue de littérature sur la taxe sur la valeur foncière.
% Même composition que le cours (cours/refaire-l-impot.tex) ; figures « nues » produites par
%   FIG_NU=1 python lvt/graphiques_lvt.py
\usepackage[T1]{fontenc}
\usepackage[utf8]{inputenc}
\usepackage{lmodern}
\usepackage[french]{babel}
\usepackage[a4paper,left=2.8cm,right=2.8cm,top=2.6cm,bottom=2.8cm]{geometry}
\usepackage{graphicx}
\usepackage[export]{adjustbox}
\usepackage{booktabs}
\usepackage{tabularx}
\usepackage{array}
\usepackage{xcolor}
\usepackage{titlesec}
\usepackage[labelfont=bf,labelsep=period,font=small,singlelinecheck=false,
            justification=raggedright,skip=6pt]{caption}
\usepackage{fancyhdr}
\usepackage{enumitem}
\usepackage{microtype}
\usepackage{textcomp}
\usepackage{amsmath}
\providecommand{\euro}{\texteuro}
\usepackage[hidelinks]{hyperref}

\definecolor{encre}{HTML}{1A1A19}
\definecolor{bleu}{HTML}{003399}
\definecolor{vert}{HTML}{1F7A4D}
\definecolor{ocre}{HTML}{C79100}
\definecolor{brique}{HTML}{993333}
\definecolor{filet}{HTML}{C8C6BE}

\titleformat{\section}{\normalfont\large\bfseries\color{bleu}}{\thesection.}{0.6em}{}
\titleformat{\subsection}{\normalfont\normalsize\bfseries\color{bleu}}{\thesubsection}{0.6em}{}
\titlespacing*{\section}{0pt}{2.4ex plus 1ex minus .2ex}{1.2ex plus .2ex}
\titlespacing*{\subsection}{0pt}{2ex plus 1ex minus .2ex}{0.8ex plus .2ex}

\pagestyle{fancy}
\fancyhf{}
\renewcommand{\headrulewidth}{0.4pt}
\fancyhead[R]{\small\thepage}

\setlength{\parindent}{0pt}
\setlength{\parskip}{0.62em}
\linespread{1.05}
\renewcommand{\arraystretch}{1.2}
\setlist{itemsep=2pt,topsep=2pt,parsep=0pt,leftmargin=1.4em}

\newenvironment{encadre}[1]
  {\par\medskip\noindent\begin{minipage}{\linewidth}%
   \color{encre}\setlength{\parskip}{0.5em}%
   \hrule height 0.8pt width \linewidth \vspace{6pt}%
   {\footnotesize\bfseries\color{vert}\MakeUppercase{#1}}\par\vspace{-2pt}%
   \small}
  {\par\vspace{4pt}\hrule height 0.4pt width \linewidth\end{minipage}\par\medskip}

\newcommand{\figbloc}[3]{%
  \begin{figure}[htbp]\centering
  \caption{#2}\label{#3}
  \includegraphics[max width=\linewidth,max height=0.40\textheight]{../figures/nu/#1.pdf}
  \IfFileExists{../figures/nu/#1_note.tex}{\par\vspace{4pt}%
    \parbox{\linewidth}{\raggedright\scriptsize\color{encre}\setlength{\parskip}{2pt}%
    \input{../figures/nu/#1_note.tex}}}{}%
  \end{figure}}

\newcommand{\pc}{\,\%}

\renewcommand{\figbloc}[3]{%
  \begin{figure}[!htbp]\centering
  \caption{#2}\label{#3}
  \includegraphics[max width=\linewidth,max height=0.30\textheight]{../figures/nu/#1.pdf}
  \par\vspace{3pt}\parbox{\linewidth}{\raggedright\scriptsize\color{encre}\setlength{\parskip}{2pt}%
    \input{../figures/nu/#1_note.tex}}%
  \end{figure}}
\usepackage{needspace}
\usepackage{etoolbox}
\preto\section{\needspace{6\baselineskip}}
\renewcommand{\topfraction}{0.9}\renewcommand{\bottomfraction}{0.8}
\renewcommand{\textfraction}{0.08}\renewcommand{\floatpagefraction}{0.85}
\setcounter{topnumber}{3}\setcounter{totalnumber}{5}
\fancyhead[L]{\small\itshape Moins taxer les ventes, mieux taxer le sol}

\begin{document}
\renewcommand{\figurename}{Figure}\renewcommand{\tablename}{Tableau}

\thispagestyle{empty}
\vspace*{0.6cm}
{\footnotesize\color{ocre}\bfseries NOTE DE TRAVAIL\quad·\quad NOVEMBRE 2026\par}
\vspace{8pt}
{\fontsize{22}{26}\selectfont\bfseries\color{bleu}Moins taxer les ventes, mieux taxer le sol\par}
\vspace{6pt}
{\large Réformer la fiscalité immobilière : diviser par deux les droits de mutation et passer à une taxe
foncière à deux taux\par}
\vspace{8pt}
{\small Colin Baget. Calculs sur le fichier de fiscalité locale 2024 (REI, DGFiP), les ventes DVF 2023-2024,
le recensement 2022, les comptes de patrimoine et l'Insee Analyses n°~91 ; programmes
\texttt{fiscalite/foncier/}. S'appuie sur la note « Taxer le terrain » (octobre 2026).\par}
\vspace{6pt}
{\color{bleu}\hrule height 1.2pt}
\vspace{4pt}

\section*{Introduction}

La France taxe beaucoup l'immobilier, et elle le taxe au mauvais endroit. Elle prélève lourdement au
moment où un logement change de mains : un ménage qui achète un logement ancien paie 5,8 à 6,3~\% du prix en
droits de mutation, en plus du prix, à l'instant précis où son apport est le plus contraint. Elle prélève
en revanche peu, et mal, sur la détention : la taxe foncière repose sur des valeurs locatives de 1970, qui
font payer proportionnellement plus les logements des communes pauvres que ceux des communes chères.

L'économie de ce choix est défavorable. Un impôt sur les ventes décourage les ventes : il retient dans leur
logement des ménages qui voudraient en changer, freine les déménagements vers un logement plus grand quand
la famille s'agrandit, et pèse d'abord sur ceux qui achètent, souvent jeunes. Un impôt récurrent sur la
valeur du sol, à l'inverse, ne décourage ni le travail ni la construction : le terrain existe de toute
façon. Il a un défaut, qui explique qu'on s'en serve si peu : il se voit, chaque année, sur un avis
d'imposition.

Cette note dresse l'état des lieux de tous les prélèvements assis sur l'immobilier, résume ce que la
recherche établit sur les droits de mutation et sur la taxe foncière, puis simule une bascule : diviser
par deux les droits de mutation, et reporter la recette sur une taxe foncière refondée, assise sur la
valeur de marché, avec un taux plus élevé sur le terrain que sur le bâti. L'objectif est une fiscalité
favorable à la croissance et à l'accession à la propriété, en particulier des jeunes ménages, sans
déstabiliser brutalement les budgets des propriétaires en place.

\section{État des lieux : 116 milliards, dont un quart sur les ventes}

Les prélèvements assis sur l'immobilier atteignent environ 116~Md€ en 2024, 4~\% du PIB
(figure~\ref{fig:inv}). La détention en fournit 69~Md€ : la taxe foncière sur les propriétés bâties
(42~Md€, 43,1 en 2025, dont environ 31,6~Md€ sur les logements), la taxe d'enlèvement des ordures ménagères
qui lui est adossée (9~Md€), la cotisation foncière des entreprises (8~Md€), l'impôt sur la fortune
immobilière (2,7~Md€) et les taxes d'habitation qui subsistent sur les résidences secondaires et les
logements vacants. Les transactions en fournissent 25~Md€ : droits de mutation (14,1~Md€ en 2024,
16,7 en 2025), TVA sur les logements neufs et les terrains à bâtir (7,1~Md€), impôt sur les plus-values
immobilières (2,8~Md€), contribution de sécurité immobilière (0,7~Md€). S'y ajoutent les droits de
succession et de donation sur l'immobilier (environ 12,5~Md€, estimation faute de ventilation
officielle), l'impôt sur les revenus locatifs (8,4~Md€) et la taxe d'aménagement (1,5~Md€).

\figbloc{fon_03_inventaire}{Les prélèvements assis sur l'immobilier en 2024}{fig:inv}

Selon l'OCDE, la France prélève en impôts périodiques sur l'immobilier 1,95~\% du PIB, le double de la
moyenne des pays membres (0,95~\%), et en impôts sur les transactions 0,62~\% contre 0,37~\%. Elle se
distingue donc des deux côtés, mais c'est sur les transactions que l'écart relatif est le plus coûteux,
parce que c'est là que l'impôt détruit le plus d'échanges par euro collecté.

\textbf{Les droits de mutation.} Pour un logement ancien, ils cumulent la part départementale (4,50~\%),
la taxe communale (1,20~\%) et les frais d'assiette de l'État (2,37~\% de la part départementale), soit
5,81~\% du prix. La loi de finances pour 2025 a permis aux départements de porter leur part à 5~\% pour les
ventes conclues du 1\textsuperscript{er}~avril~2025 au 31~mars~2028 ; 88 collectivités sur une centaine
l'ont fait, ce qui porte le taux à 6,32~\%. Les primo-accédants achetant leur résidence principale sont
exclus de cette majoration ; la possibilité d'une réduction propre aux primo-accédants n'est utilisée que
par la Savoie.

\textbf{La taxe foncière.} Elle repose sur la moitié de la valeur locative cadastrale, établie selon les loyers
de 1970 et revalorisée depuis par des coefficients nationaux. La révision des valeurs des logements a été
reportée à 2031 par la loi de finances pour 2026. Rapportée à la valeur de marché, elle pèse 0,55~\% dans le
dixième des logements situés dans les communes les moins chères et 0,16~\% dans le dixième le plus cher
(figure~\ref{fig:tfv}) ; à Paris, 0,10~\%. L'Insee aboutit au même constat par une autre voie : la taxe
foncière est régressive rapportée au patrimoine immobilier et progressive rapportée au revenu (André et
Meslin, 2023).

\section{Ce que dit la recherche}

\subsection*{Les droits de mutation réduisent les ventes et la mobilité}

Trois résultats sont solides. D'abord, les ventes réagissent fortement au calendrier : le mois qui précède
une hausse, elles bondissent de 26~\% en France en 2014 (Bérard et Trannoy, 2018), de 41~\% en Allemagne ;
ce sont des ventes avancées, pas des ventes créées. Ensuite, au-delà de ces effets de calendrier, la plupart
des études trouvent que chaque point de droits supprime durablement environ 8~\% des ventes, dans une
fourchette de 0 à 18~\% (Dachis, Duranton et Turner, 2012, à Toronto ; Fritzsche et Vandrei, 2019, en
Allemagne ; Best et Kleven, 2018, au Royaume-Uni). Enfin, la mobilité des propriétaires baisse de 8 à 14~\%
par point (Pays-Bas, Finlande, Toronto). En Finlande, ce sont surtout les déménagements courts et les
déménagements vers plus grand qui disparaissent, pas les réductions de taille (Eerola et al., 2021) : la taxe
frappe précisément le ménage qui a besoin d'une pièce de plus parce qu'un enfant arrive.

Pour la France, les preuves sont plus minces et partagées. Après la hausse de 0,7~point de 2014, Bérard et
Trannoy trouvent une baisse des ventes de 6~\% sur trois mois, puis plus rien ; Poulhès (2017) trouve une
baisse durable de 10 à 12~\%, de 20~\% dans les marchés détendus et nulle dans les marchés tendus. Il trouve
aussi que le prix hors droits n'a pas bougé : l'acheteur a supporté toute la hausse. Une baisse des droits
profiterait donc d'abord, au moins à court terme, aux acheteurs ; ailleurs, les baisses durables passent en
partie dans les prix au profit des vendeurs, et une exonération ciblée sur les primo-accédants au
Royaume-Uni a augmenté leurs achats de 11~\% tout en relevant de 1~\% les prix qu'ils payaient (HMRC, 2023).

Le coût économique des droits de mutation croît vite avec leur taux, comme celui de tout impôt : autour de
1 à 2~\%, il représente 5 à 30~\% de la recette ; vers 4 à 5~\%, il approche 0,7 à 0,9~€ par euro
supplémentaire collecté (Buettner, en Allemagne ; Määttänen et Terviö, 2022). Au taux français, il serait
supérieur à 2~€ par euro selon le modèle finlandais, mais personne ne l'a estimé pour la France. Le seul
modèle français (Bontemps, Cherbonnier et Magnac, 2023) trouve un gain modeste à remplacer les droits par
un impôt foncier (+0,2~\% de consommation), qui s'inverse si l'on finance la baisse par l'impôt sur le
revenu : c'est la raison de la reporter sur la détention.

\subsection*{La taxe foncière : un bon impôt mal assis}

Le principal inconvénient d'un impôt récurrent sur l'immobilier est la transition. Une hausse permanente se
capitalise dans les prix : le propriétaire présent au moment de l'annonce perd la valeur actuelle de toutes
les hausses futures, et les acheteurs suivants achètent moins cher (Koster et Pinchbeck, 2022, pour la
council tax anglaise). C'est un transfert des propriétaires en place vers les futurs acheteurs. Pour une
taxe assise sur le seul sol, l'ampleur de cette capitalisation reste discutée (Danemark : complète selon
Høj et al., 2018, nulle selon Nielsson et al., 2024). Un ménage âgé riche en patrimoine mais modeste en
revenu peut aussi avoir du mal à payer ; les reports de paiement facultatifs sont peu utilisés, sauf quand
ils sont proposés par défaut, comme au Danemark.

L'assiette est l'autre difficulté. Les évaluations de masse sont souvent régressives, parce qu'elles suivent
mal la valeur des quartiers chers (Berry, 2021 ; Avenancio-León et Howard, 2022). Philadelphie a corrigé
trente ans de gel par une réévaluation générale en 2014, accompagnée d'un abattement et d'un dispositif
pour les occupants de longue date. Pittsburgh, à l'inverse, a abandonné en 2001 son double taux après une
réévaluation brutale qui a fait bondir la valeur du sol de 81~\% sans lissage. Un plafond de hausse attaché
au propriétaire plutôt qu'au bien crée un verrou : en Californie, la Proposition~13 retient les ménages dans
leur logement (Ferreira, 2010).

Le double taux, plus élevé sur le terrain que sur le bâti, a un effet modeste mais positif sur la
construction : en Pennsylvanie, il a accru de 5 à 6~points la croissance décennale du nombre de pièces
(Banzhaf et Lavery, 2010). La part de l'impôt qui porte sur le bâti décourage l'investissement là où l'offre
de logements peut réagir (Lutz, 2015). Le Bade-Wurtemberg est allé au bout de la logique : depuis 2025, sa
taxe foncière ne porte que sur la valeur du sol, établie par des valeurs de zone, avec un taux réduit de
30~\% pour le logement ; la plus haute juridiction fiscale allemande l'a jugée conforme en mai 2026. À
Karlsruhe, à recette constante, les maisons individuelles paient 76~\% de plus, les copropriétés 20~\% de
moins et les terrains nus 265~\% de plus : l'assiette sol charge les grandes parcelles et allège la ville
dense.

\section{La proposition et ce qu'elle donne}

\begin{encadre}{Comment la bascule est chiffrée}
\textbf{Champ.} Les 30,5~millions de logements des particuliers (parc du recensement moins les HLM) de France
métropolitaine hors Alsace-Moselle, où les ventes ne sont pas publiées. Leur valeur est de 6~760~Md€ : prix
médian au m² des ventes 2023-2024 de la commune (sinon de l'intercommunalité ou du département), multiplié
par la surface médiane des logements vendus.

\textbf{Part du terrain.} Pour les maisons, prix au m² des terrains à bâtir multiplié par la surface de la
parcelle ; quand la commune compte trop peu de ventes de terrains, le prix départemental est ajusté selon
l'écart de prix des maisons (élasticité estimée de 1,54 sur 3~300 communes). Pour les appartements, valeur
du logement moins le coût de reconstruction du bâti, calé pour que l'ensemble retrouve la part du terrain
des comptes de patrimoine des ménages (45~\%). Elle va de 16~\% dans le dixième le moins cher à 67~\% dans le
plus cher.

\textbf{Charges actuelles.} Taxe foncière : rapport taxe/valeur de chaque commune selon l'Insee (2017),
appliqué aux valeurs de 2024 et calé sur le total (23,6~Md€ pour ce champ). Droits de mutation : 5,81~\% des
ventes de logements anciens, 11,4~Md€ par an en moyenne en 2023-2024.

\textbf{Ce qui n'est pas modélisé.} Les effets sur les prix (capitalisation), l'âge et le revenu des
propriétaires, les locaux professionnels, les logements sociaux et les résidences secondaires en tant que
telles.
\end{encadre}

\textbf{Première étape : diviser par deux les droits de mutation.} Le taux passe de 5,81 à 2,9~\%. La perte de
recette est de 5,7~Md€ si les ventes ne réagissent pas, de 4,6~Md€ si chaque point de baisse augmente les ventes
de 6~\%, valeur centrale de la littérature (4,3 à 5,7~Md€ selon l'hypothèse). Une baisse de 30~\% coûterait 2,6
à 3,4~Md€. La baisse doit s'appliquer dès son annonce, sans délai : sinon les ventes s'arrêtent jusqu'à son
entrée en vigueur. La Finlande l'a fait en 2023 en la rendant rétroactive à la date du projet de loi.

\textbf{Deuxième étape : reporter la recette sur la taxe foncière.} Trois manières de le faire ont des
effets très différents (figures~\ref{fig:tfv} et~\ref{fig:var}).

La plus simple relève la taxe foncière actuelle de 24~\% partout (20~\% si les ventes réagissent) : entre 110
et 230~€ de plus par logement et par an. Elle ne bouleverse rien, mais elle aggrave l'injustice des
valeurs de 1970 : la taxe atteindrait 0,68~\% de la valeur dans les communes les moins chères contre 0,20~\%
dans les plus chères.

La plus ambitieuse remplace la taxe foncière par un impôt national à deux taux sur la valeur de marché :
0,73~\% sur le terrain et 0,18~\% sur le bâti. Elle supprime les distorsions de 1970, mais d'un coup : la taxe
augmenterait de 238~\% en moyenne dans le dixième des logements le plus cher et un logement sur quatre
verrait sa taxe augmenter de plus de moitié. Sans une transition de plusieurs décennies, c'est impossible.

\figbloc{fon_01_tf_valeur}{La taxe foncière rapportée à la valeur des logements, aujourd'hui et après la bascule}{fig:tfv}

\textbf{Notre proposition} combine les deux logiques. Chaque commune garde sa recette de taxe foncière et son
pouvoir de taux, mais l'assiette devient la valeur de marché, avec un taux sur le terrain égal à quatre fois
le taux sur le bâti, dans la fourchette pratiquée en Pennsylvanie. La perte de droits de mutation est
remplacée par un taux national unique sur la valeur des terrains, 0,15~\% (0,19~\% si les ventes ne
réagissent pas), reversé aux départements selon leurs pertes. Le taux national évite une difficulté
sérieuse : si chaque département remplaçait seul ses droits perdus, il lui faudrait un taux de 0,12~\% de la
valeur du sol à Paris mais de 0,85~\% dans la Creuse, où les terrains valent peu.

\figbloc{fon_02_variation}{Variation de la charge par logement et par an, selon le prix au m² de la commune}{fig:var}

Avec la réaction centrale des ventes, la taxe foncière augmente de 2~\% dans le dixième des logements le moins
cher et de 60~\% dans le plus cher, de 121~€ en moyenne pour une maison et de 195~€ pour un appartement. En
comptant la baisse des droits de mutation, rapportée à l'ensemble des logements de la commune, la charge
totale varie de $-$86 à +26~€ par logement et par an selon le décile. Un logement sur six (18~\%) voit sa
taxe foncière augmenter de plus de 300~€. Rapportée à la valeur, elle va de 0,56~\% dans les communes les
moins chères à 0,26~\% dans les plus chères : l'écart actuel se resserre sans être renversé.

\figbloc{fon_04_primo}{Ce que gagne un ménage qui achète le logement médian de sa commune}{fig:primo}

\textbf{Pour les jeunes acheteurs, le gain est immédiat} (figure~\ref{fig:primo}). Un ménage qui achète le
logement médian de sa commune paie entre 2~300~€ (un appartement à Saint-Étienne) et 22~700~€ (une maison à
Arcachon) de droits en moins, au moment où son apport est le plus contraint : 8~100~€ pour un appartement à
Lyon, 13~000~€ pour une maison à Nantes, 12~100~€ pour un appartement à Paris. La hausse de taxe foncière qu'il
paie ensuite (de $-$50 à +660~€ par an selon la commune, dans l'hypothèse prudente où les ventes ne
réagissent pas) ne rattrape cette économie qu'au bout de 17~ans au moins (Bordeaux), et de plusieurs décennies
dans la plupart des cas. Le perdant est le propriétaire installé qui ne vend pas : il paie la hausse sans
profiter de la baisse, jusqu'à 22~000~€ en valeur actuelle pour un appartement parisien. Si la hausse est capitalisée, cette perte se reporte sur le prix de son bien au moment de
l'annonce, et le prix que paie le jeune acheteur baisse d'autant : la bascule transfère de la valeur des
propriétaires en place, souvent âgés, vers les acheteurs, souvent jeunes.

\textbf{Pour les familles}, le gain tient au type de déménagement que les droits découragent le plus :
l'agrandissement. Avec des droits divisés par deux, passer d'un trois-pièces à un quatre-pièces à la
naissance d'un deuxième enfant coûte la moitié moins en impôt. La proposition charge aussi un peu moins les
maisons que les appartements, parce que le terrain des maisons hors des grandes villes vaut peu. Aucune étude
ne relie directement les droits de mutation à la natalité : l'effet attendu passe par le logement, pas par
une incitation directe.

\textbf{Pour les ménages les plus aisés}, l'effet net est limité. La taxe foncière augmente davantage dans
les communes chères (en moyenne 450~€ par an dans le dixième le plus cher, avec la réaction des ventes), mais
les acheteurs de biens chers économisent davantage de droits. L'impôt sur la fortune immobilière n'est pas
modifié.

\section{Comment passer d'un système à l'autre}

\textbf{Construire l'assiette.} La France n'a pas de valeurs du sol par zone, comme les
\emph{Bodenrichtwerte} allemandes. Les ventes publiées (DVF) et leur version enrichie réservée aux acteurs
publics (DV3F) permettent de les établir par commune ou par quartier, en combinant les ventes de terrains et
les prix des logements, comme le fait cette note. L'administration fiscale établirait des valeurs de zone,
révisées tous les trois ans et lissées sur trois ans ; le propriétaire pourrait contester par expertise
quand sa valeur s'écarte nettement du marché, comme en Bade-Wurtemberg au-delà de 30~\%. Il faut écarter
la révision de la valeur à chaque vente, adoptée en commission dans le projet de loi de finances pour 2027
(amendement I-CF1872) : elle créerait le même verrou que la Proposition~13.

\textbf{Échelonner la hausse, pas la baisse.} Les droits de mutation baissent en une fois, dès l'annonce. La
nouvelle assiette entre en vigueur sur dix ans, avec un plafond de hausse annuelle attaché au bien et non au
propriétaire (le Danemark plafonne la hausse à 4,75~\% par an). Les ménages âgés aux revenus modestes, déjà
exonérés au-delà de 75~ans sous condition de ressources, se verraient proposer par défaut un report de
paiement remboursé à la vente ou à la succession.

\textbf{Garder l'autonomie locale.} Les communes votent leur taux sur la nouvelle assiette ; la loi fixe le
rapport entre le taux sur le terrain et le taux sur le bâti. Le Conseil constitutionnel n'exige pas de
compensation à l'euro près quand une recette locale est supprimée (décision 2019-796 DC) ; le taux national
sur le terrain, reversé aux départements, assure néanmoins cette compensation.

\textbf{Limites.} Les valeurs sont des médianes communales, pas des valeurs de biens ; la part du terrain
des appartements est un calage sur les comptes nationaux, que l'Insee obtient lui-même par différence ; la
charge actuelle de chaque commune repose sur un rapport de 2017. L'effet de la taxe foncière française sur
les loyers et la construction n'a jamais été mesuré, et la répartition par âge de la bascule reste à
établir sur des données individuelles.

\vspace{6pt}
{\scriptsize\textbf{Références.} André et Meslin (2023), \emph{Insee Analyses} n°~91. Avenancio-León et Howard
(2022), « The Assessment Gap », \emph{QJE} 137(3). Banzhaf et Lavery (2010), « Can the Land Tax Help Curb Urban
Sprawl? », \emph{Journal of Urban Economics} 67(2). Bérard et Trannoy (2018), \emph{Économie et Statistique}
n°~500-501-502. Berry (2021), « Reassessing the Property Tax », Harris School. Best et Kleven (2018),
\emph{Review of Economic Studies} 85(1). Bontemps, Cherbonnier et Magnac (2023), « Reducing Transaction Taxes on
Housing in Highly Regulated Economies », TSE WP 23-14. Buettner (2017), CESifo WP 6321. Dachis, Duranton et
Turner (2012), \emph{Journal of Economic Geography} 12(2). Eerola, Harjunen, Lyytikäinen et Saarimaa (2021),
« Revisiting the Effects of Housing Transfer Taxes », \emph{Journal of Urban Economics} 124. Ferreira (2010),
\emph{Journal of Public Economics} 94(9-10). Fritzsche et Vandrei (2019), \emph{Regional Science and Urban
Economics} 74. HMRC (2023), évaluation de l'allègement pour primo-accédants. Høj, Jørgensen et Schou (2018),
\emph{Fiscal Studies} 39(2). Koster et Pinchbeck (2022), \emph{American Economic Journal: Economic Policy} 14(1).
Lutz (2015), \emph{American Economic Journal: Economic Policy} 7(1). Määttänen et Terviö (2022), \emph{Economic
Journal} 132(644). Nielsson, Wroblewski et Yding (2024), document de travail. Poulhès (2017), \emph{Théma},
CGDD. OCDE, \emph{Revenue Statistics}. DGCL-DGFiP, note d'analyse du REI 2025. SDES, rapport du compte du
logement 2025. Conseil constitutionnel, décision 2019-796 DC. Détail des sources et statut de lecture :
\texttt{foncier/notes/}.\par}

\end{document}
